\documentclass[12pt,a4paper]{article}

\usepackage[T2A]{fontenc}
\usepackage[utf8]{inputenc}
\usepackage[russian]{babel}
\usepackage{amsmath,amssymb,amsthm}
\usepackage{array}
\usepackage{booktabs}
\usepackage{enumitem}
\usepackage{geometry}
\usepackage{longtable}
\usepackage{hyperref}

\geometry{left=25mm,right=20mm,top=20mm,bottom=25mm}

\newtheorem{definition}{Определение}
\newtheorem{remark}{Замечание}
\newtheorem{example}{Пример}
\newtheorem{principle}{Принцип}

\title{Билет №61 \\ \small Требования к программному продукту (надежность, переносимость, познаваемость, рациональная ресурсоемкость) и их влияние на системы программирования и технологии разработки программных систем. (СпБПУ)}
\author{Сериков Владислав Александрович}
\date{16 мая 2026}

\begin{document}

\maketitle
\tableofcontents
\newpage

\section{Понятие программного продукта и требований к нему}

\begin{definition}
\textbf{Программный продукт} --- это совокупность программ, данных, документации, конфигурационных файлов, средств установки, сопровождения и других артефактов, предназначенных для поставки пользователю или эксплуатации в заданной среде.
\end{definition}

Программный продукт рассматривается не только как исходный код или исполняемая программа. В его состав обычно входят:
\begin{itemize}
    \item исполняемые модули;
    \item исходные тексты;
    \item библиотеки и внешние зависимости;
    \item пользовательская и техническая документация;
    \item тесты;
    \item скрипты сборки, установки и развертывания;
    \item средства мониторинга и диагностики;
    \item спецификации интерфейсов;
    \item эксплуатационные инструкции.
\end{itemize}

\begin{definition}
\textbf{Требование к программному продукту} --- это зафиксированное свойство, функция, ограничение или условие, которому должен удовлетворять программный продукт для решения задач заинтересованных сторон.
\end{definition}

Требования делят на две большие группы:
\begin{enumerate}
    \item \textbf{Функциональные требования} --- описывают, \emph{что} должна делать система.
    \item \textbf{Нефункциональные требования} --- описывают, \emph{как хорошо} система должна выполнять свои функции и в каких условиях она должна работать.
\end{enumerate}

К нефункциональным требованиям относятся надежность, переносимость, познаваемость, рациональная ресурсоемкость, безопасность, сопровождаемость, совместимость и другие характеристики качества.

\section{Качество программного продукта}

\begin{definition}
\textbf{Качество программного продукта} --- это совокупность свойств программного продукта, определяющих его способность удовлетворять заданные и предполагаемые потребности пользователей при заданных условиях эксплуатации.
\end{definition}

Качество программного продукта является многомерным понятием. Один и тот же продукт может быть очень быстрым, но плохо переносимым; надежным, но сложным для понимания; удобным для пользователя, но трудно сопровождаемым.

В инженерии программного обеспечения для описания качества часто используют модели качества. В современных стандартах, например в семействе ISO/IEC 25000 SQuaRE, выделяются такие характеристики, как:
\begin{itemize}
    \item функциональная пригодность;
    \item производительная эффективность;
    \item совместимость;
    \item удобство использования;
    \item надежность;
    \item защищенность;
    \item сопровождаемость;
    \item переносимость.
\end{itemize}

В данном конспекте подробно рассматриваются четыре важные характеристики:
\begin{enumerate}
    \item надежность;
    \item переносимость;
    \item познаваемость;
    \item рациональная ресурсоемкость.
\end{enumerate}

\section{Надежность программного продукта}

\subsection{Определение надежности}

\begin{definition}
\textbf{Надежность программного продукта} --- это способность программного продукта сохранять требуемый уровень функционирования в заданных условиях в течение заданного времени.
\end{definition}

Иначе говоря, надежная программа:
\begin{itemize}
    \item корректно выполняет свои функции;
    \item устойчиво работает при допустимых входных данных;
    \item обнаруживает и обрабатывает ошибки;
    \item сохраняет данные при сбоях;
    \item восстанавливает работоспособность после отказов.
\end{itemize}

Надежность тесно связана с понятиями ошибки, дефекта, отказа и сбоя.

\begin{definition}
\textbf{Ошибка} --- неправильное действие человека, допущенное при проектировании, программировании, настройке или эксплуатации системы.
\end{definition}

\begin{definition}
\textbf{Дефект} --- наличие в программном артефакте неправильного или неполного элемента, который может привести к неправильному поведению программы.
\end{definition}

\begin{definition}
\textbf{Отказ} --- наблюдаемое отклонение поведения системы от требуемого поведения.
\end{definition}

\begin{definition}
\textbf{Сбой} --- событие, приводящее к нарушению нормального функционирования системы.
\end{definition}

Связь этих понятий можно представить следующим образом:
\[
\text{ошибка человека} \rightarrow \text{дефект в продукте} \rightarrow \text{отказ при выполнении}.
\]

\subsection{Составляющие надежности}

Надежность программного продукта обычно включает следующие подхарактеристики:

\begin{enumerate}
    \item \textbf{Безотказность} --- способность работать без отказов в течение заданного времени.
    \item \textbf{Отказоустойчивость} --- способность продолжать работу при наличии отдельных ошибок или отказов компонентов.
    \item \textbf{Восстанавливаемость} --- способность возвращаться в работоспособное состояние после отказа.
    \item \textbf{Доступность} --- способность быть готовым к использованию в нужный момент времени.
    \item \textbf{Целостность данных} --- способность не допускать потери или некорректного изменения данных.
\end{enumerate}

\subsection{Количественные показатели надежности}

Для оценки надежности используют метрики.

\subsubsection{Вероятность безотказной работы}

Если $T$ --- случайная величина, обозначающая время до отказа, то вероятность безотказной работы на интервале $[0,t]$ определяется как:
\[
R(t) = P(T > t).
\]

Здесь:
\begin{itemize}
    \item $R(t)$ --- функция надежности;
    \item $T$ --- время до отказа;
    \item $t$ --- рассматриваемый момент времени.
\end{itemize}

\subsubsection{Интенсивность отказов}

Интенсивность отказов показывает, как часто происходят отказы в единицу времени:
\[
\lambda = \frac{N}{T_{\Sigma}},
\]
где:
\begin{itemize}
    \item $N$ --- число отказов;
    \item $T_{\Sigma}$ --- суммарное время работы системы.
\end{itemize}

\subsubsection{Среднее время наработки на отказ}

\begin{definition}
\textbf{MTTF} от английского \emph{Mean Time To Failure} --- среднее время до отказа для невосстанавливаемой системы.
\end{definition}

\[
MTTF = \frac{1}{\lambda}.
\]

\begin{definition}
\textbf{MTBF} от английского \emph{Mean Time Between Failures} --- среднее время между отказами для восстанавливаемой системы.
\end{definition}

\subsubsection{Среднее время восстановления}

\begin{definition}
\textbf{MTTR} от английского \emph{Mean Time To Repair} --- среднее время восстановления системы после отказа.
\end{definition}

\subsubsection{Доступность}

Доступность можно оценить формулой:
\[
A = \frac{MTBF}{MTBF + MTTR}.
\]

Если $MTBF = 1000$ часов, а $MTTR = 1$ час, то:
\[
A = \frac{1000}{1000 + 1} \approx 0.999.
\]

Это означает доступность примерно $99.9\%$.

\subsection{Методы повышения надежности}

Основные методы повышения надежности программного продукта:

\begin{enumerate}
    \item \textbf{Формализация требований.}

    Чем точнее описаны требования, тем меньше вероятность дефектов, связанных с неправильным пониманием задачи.

    \item \textbf{Модульность.}

    Система разбивается на относительно независимые компоненты. Это уменьшает распространение ошибок.

    \item \textbf{Инкапсуляция.}

    Внутреннее состояние модуля скрывается, доступ к нему осуществляется через контролируемые интерфейсы.

    \item \textbf{Проверка входных данных.}

    Программа не должна предполагать, что входные данные всегда корректны.

    \item \textbf{Исключения и обработка ошибок.}

    Ошибки должны обнаруживаться и обрабатываться явно.

    \item \textbf{Тестирование.}

    Используются модульные, интеграционные, системные, регрессионные, нагрузочные и приемочные тесты.

    \item \textbf{Статический анализ.}

    Исходный код проверяется без запуска программы.

    \item \textbf{Динамический анализ.}

    Программа проверяется во время выполнения.

    \item \textbf{Резервирование.}

    Для критических компонентов создаются резервные экземпляры.

    \item \textbf{Журналирование.}

    Система фиксирует важные события, ошибки и действия пользователей.

    \item \textbf{Мониторинг.}

    В рабочей среде контролируются показатели состояния системы.

    \item \textbf{Автоматическое восстановление.}

    При отказах система может перезапускать компоненты, переключаться на резервные узлы или восстанавливать состояние из контрольной точки.
\end{enumerate}

\subsection{Пример проектного решения для повышения надежности}

Пусть разрабатывается сервис обработки платежей. Для повышения надежности используются:
\begin{itemize}
    \item транзакции базы данных;
    \item повторные попытки при временных сетевых ошибках;
    \item идемпотентные операции;
    \item журналирование всех операций;
    \item резервное копирование;
    \item мониторинг доступности;
    \item автоматические тесты.
\end{itemize}

Минимальный пример идемпотентной операции:

\begin{verbatim}
POST /payments
Idempotency-Key: abc-123

{
  "amount": 1000,
  "currency": "RUB"
}
\end{verbatim}

Если клиент повторит запрос с тем же ключом \texttt{abc-123}, сервер не должен создать второй платеж, а должен вернуть результат первой операции.

\subsection{Влияние требований надежности на системы программирования}

Требования надежности влияют на языки программирования, библиотеки, среды разработки и инструменты следующим образом:

\begin{enumerate}
    \item \textbf{Строгая типизация.}

    Языки со строгой статической типизацией позволяют обнаруживать часть ошибок до запуска программы.

    \item \textbf{Механизмы исключений.}

    Язык должен поддерживать структурированную обработку исключительных ситуаций.

    \item \textbf{Управление памятью.}

    Автоматическое управление памятью снижает риск утечек памяти и ошибок обращения к освобожденной памяти.

    \item \textbf{Контракты и утверждения.}

    Используются предусловия, постусловия и инварианты.

    \item \textbf{Средства модульного тестирования.}

    Современные системы программирования включают фреймворки тестирования.

    \item \textbf{Статические анализаторы.}

    Важны инструменты, обнаруживающие потенциальные ошибки в коде.

    \item \textbf{Отладчики и профилировщики.}

    Помогают находить причины отказов и узкие места.

    \item \textbf{Средства непрерывной интеграции.}

    Позволяют автоматически запускать тесты при каждом изменении кода.

    \item \textbf{Средства контроля версий.}

    Позволяют отслеживать изменения и быстро откатываться к стабильным версиям.
\end{enumerate}

\subsection{Влияние требований надежности на технологию разработки}

При высоких требованиях надежности технология разработки становится более строгой:
\begin{itemize}
    \item требования документируются формально;
    \item архитектура проходит экспертную проверку;
    \item код проверяется рецензированием;
    \item тестирование становится обязательной частью процесса;
    \item изменения проходят через систему контроля качества;
    \item применяется регрессионное тестирование;
    \item используется управление конфигурациями;
    \item вводятся процедуры анализа отказов.
\end{itemize}

Для критических систем могут применяться:
\begin{itemize}
    \item формальные методы;
    \item математическое доказательство свойств программ;
    \item моделирование;
    \item верификация моделей;
    \item сертификация по отраслевым стандартам.
\end{itemize}

\section{Переносимость программного продукта}

\subsection{Определение переносимости}

\begin{definition}
\textbf{Переносимость программного продукта} --- это способность программного продукта быть перенесенным из одной аппаратной, программной, организационной или эксплуатационной среды в другую с приемлемыми затратами.
\end{definition}

Переносимая программа может работать:
\begin{itemize}
    \item на разных операционных системах;
    \item на разных аппаратных архитектурах;
    \item с разными системами управления базами данных;
    \item в разных облачных средах;
    \item в разных локалях и языковых окружениях.
\end{itemize}

\subsection{Составляющие переносимости}

Переносимость включает:

\begin{enumerate}
    \item \textbf{Адаптируемость} --- возможность приспособления к новой среде.
    \item \textbf{Устанавливаемость} --- возможность установки и удаления в заданной среде.
    \item \textbf{Заменяемость} --- возможность заменить другим продуктом или заменить существующий продукт.
    \item \textbf{Совместимость с платформами} --- способность работать на разных ОС, процессорах, версиях библиотек.
    \item \textbf{Локализуемость} --- возможность адаптации к языку, культуре, формату дат, валют и чисел.
\end{enumerate}

\subsection{Источники непереносимости}

Основные причины плохой переносимости:
\begin{itemize}
    \item использование нестандартных возможностей языка;
    \item жесткая привязка к операционной системе;
    \item зависимость от абсолютных путей файловой системы;
    \item зависимость от конкретного формата кодировки;
    \item зависимость от порядка байтов;
    \item использование платформо-зависимых библиотек;
    \item смешивание бизнес-логики с кодом доступа к среде;
    \item отсутствие автоматизированной сборки;
    \item отсутствие контейнеризации или описания окружения.
\end{itemize}

\subsection{Методы повышения переносимости}

Для повышения переносимости применяются:

\begin{enumerate}
    \item \textbf{Использование стандартных языков и библиотек.}

    Например, предпочтительно использовать стандартные средства языка вместо специфичных API конкретной ОС.

    \item \textbf{Абстрагирование от платформы.}

    Платформо-зависимый код помещается в отдельный слой.

    \item \textbf{Использование виртуальных машин и интерпретаторов.}

    Примеры: JVM, .NET CLR, Python runtime.

    \item \textbf{Контейнеризация.}

    Приложение поставляется вместе с описанием окружения.

    \item \textbf{Кроссплатформенные фреймворки.}

    Используются библиотеки, поддерживающие несколько платформ.

    \item \textbf{Автоматизация сборки.}

    Сборка должна быть воспроизводимой.

    \item \textbf{Конфигурируемость.}

    Параметры среды не должны быть жестко записаны в коде.

    \item \textbf{Тестирование на нескольких платформах.}

    Нельзя считать программу переносимой, если она реально не проверялась в целевых средах.
\end{enumerate}

\subsection{Минимальный пример выделения платформо-зависимого слоя}

Пусть программа должна получать путь к домашнему каталогу пользователя.

Плохой вариант:
\begin{verbatim}
home = "C:\\Users\\Ivan"
\end{verbatim}

Такой код привязан к Windows и конкретному пользователю.

Лучший вариант:
\begin{verbatim}
home = get_user_home_directory()
\end{verbatim}

Внутри функции можно реализовать разные варианты для разных ОС:
\begin{verbatim}
function get_user_home_directory():
    if OS == "Windows":
        return getenv("USERPROFILE")
    else:
        return getenv("HOME")
\end{verbatim}

В этом случае остальная программа не зависит от деталей операционной системы.

\subsection{Влияние требований переносимости на системы программирования}

Требования переносимости приводят к появлению и развитию:

\begin{enumerate}
    \item \textbf{Стандартизованных языков программирования.}

    Чем строже определен язык, тем легче переносить программы между компиляторами и платформами.

    \item \textbf{Кроссплатформенных библиотек.}

    Например, библиотеки для работы с файлами, сетью, потоками, графическим интерфейсом.

    \item \textbf{Виртуальных машин.}

    Виртуальная машина скрывает различия между аппаратными и операционными платформами.

    \item \textbf{Промежуточных представлений.}

    Программа компилируется не сразу в машинный код, а в байт-код или промежуточный язык.

    \item \textbf{Систем сборки.}

    Они позволяют описывать процесс сборки независимо от конкретной среды.

    \item \textbf{Пакетных менеджеров.}

    Они фиксируют версии зависимостей и упрощают развертывание.

    \item \textbf{Контейнерных технологий.}

    Они позволяют описывать окружение приложения как часть поставки.
\end{enumerate}

\subsection{Влияние требований переносимости на технологию разработки}

Если переносимость является важным требованием, то в технологию разработки включают:
\begin{itemize}
    \item матрицу поддерживаемых платформ;
    \item правила запрета платформо-зависимого кода вне специальных модулей;
    \item кроссплатформенное тестирование;
    \item автоматическую сборку для разных ОС;
    \item хранение конфигурации отдельно от кода;
    \item документирование требований к окружению;
    \item регрессионные тесты на всех поддерживаемых платформах.
\end{itemize}

\section{Познаваемость программного продукта}

\subsection{Смысл понятия}

Термин \emph{познаваемость} в контексте качества программного продукта можно рассматривать в двух близких смыслах:

\begin{enumerate}
    \item \textbf{Познаваемость для пользователя} --- насколько легко пользователь может понять назначение продукта, освоить его функции и научиться работать с ним.
    \item \textbf{Познаваемость для разработчика и сопровождающего персонала} --- насколько легко понять внутреннее устройство программы, ее архитектуру, код, данные и поведение.
\end{enumerate}

Первый смысл близок к характеристикам удобства использования: понятность, обучаемость, простота освоения.

Второй смысл близок к сопровождаемости: анализируемость, читаемость, документированность, модульность.

\begin{definition}
\textbf{Познаваемость программного продукта} --- это способность программного продукта быть понятым, изученным и освоенным пользователями, разработчиками или сопровождающим персоналом с приемлемыми затратами времени и усилий.
\end{definition}

\subsection{Познаваемость для пользователя}

Пользовательская познаваемость определяется тем, насколько легко пользователь может:
\begin{itemize}
    \item понять назначение продукта;
    \item обнаружить нужные функции;
    \item научиться выполнять типовые операции;
    \item предсказать результат действия;
    \item исправить ошибку;
    \item использовать справку и документацию.
\end{itemize}

Факторы, повышающие пользовательскую познаваемость:
\begin{enumerate}
    \item единообразный интерфейс;
    \item ясные названия команд;
    \item понятные сообщения об ошибках;
    \item наличие подсказок;
    \item обучающие сценарии;
    \item документация;
    \item соответствие пользовательским ожиданиям;
    \item минимизация скрытых режимов работы.
\end{enumerate}

\subsection{Познаваемость для разработчика}

Разработческая познаваемость определяется тем, насколько легко понять:
\begin{itemize}
    \item структуру системы;
    \item назначение модулей;
    \item связи между компонентами;
    \item формат данных;
    \item алгоритмы;
    \item причины проектных решений;
    \item порядок сборки и запуска;
    \item сценарии тестирования.
\end{itemize}

Факторы, повышающие разработческую познаваемость:
\begin{enumerate}
    \item модульность;
    \item слабая связность компонентов;
    \item высокая связность внутри модуля;
    \item единый стиль кодирования;
    \item осмысленные имена;
    \item комментарии там, где они объясняют нетривиальные решения;
    \item архитектурная документация;
    \item диаграммы;
    \item тесты как исполняемая спецификация;
    \item простота сборки и запуска.
\end{enumerate}

\subsection{Познаваемость и сложность}

Познаваемость тесно связана со сложностью программы.

Сложность бывает:
\begin{itemize}
    \item \textbf{алгоритмическая};
    \item \textbf{структурная};
    \item \textbf{когнитивная};
    \item \textbf{интерфейсная};
    \item \textbf{организационная}.
\end{itemize}

Чем выше необоснованная сложность, тем ниже познаваемость.

\begin{principle}
\textbf{Принцип минимизации сложности.}
При прочих равных следует выбирать такое проектное решение, которое проще понять, проверить, изменить и объяснить.
\end{principle}

\subsection{Метрики познаваемости}

Познаваемость трудно измерить напрямую, но можно использовать косвенные показатели:
\begin{itemize}
    \item размер модуля в строках кода;
    \item цикломатическая сложность;
    \item глубина вложенности;
    \item число зависимостей;
    \item число публичных методов;
    \item покрытие документацией;
    \item время, необходимое новому разработчику для выполнения первой задачи;
    \item число ошибок пользователя при выполнении типового сценария;
    \item время обучения пользователя.
\end{itemize}

\subsection{Цикломатическая сложность}

Одной из известных метрик структурной сложности является цикломатическая сложность Маккейба.

Для программы, представленной в виде графа потока управления:
\[
V(G) = E - N + 2P,
\]
где:
\begin{itemize}
    \item $E$ --- число ребер графа;
    \item $N$ --- число вершин графа;
    \item $P$ --- число связных компонент.
\end{itemize}

Для одной функции обычно $P = 1$, поэтому:
\[
V(G) = E - N + 2.
\]

Интуитивно цикломатическая сложность показывает число линейно независимых путей выполнения. Чем она выше, тем труднее понять и протестировать функцию.

\subsection{Пример влияния сложности на познаваемость}

Плохой вариант:

\begin{verbatim}
if (a && !b || c && d || !a && e) {
    process1();
} else {
    process2();
}
\end{verbatim}

Более познаваемый вариант:

\begin{verbatim}
userHasAccess = a && !b
adminOverride = c && d
emergencyMode = !a && e

if (userHasAccess || adminOverride || emergencyMode) {
    process1()
} else {
    process2()
}
\end{verbatim}

Второй вариант может быть немного длиннее, но он легче для понимания, так как сложное условие разделено на именованные части.

\subsection{Влияние требований познаваемости на системы программирования}

Требования познаваемости способствуют развитию:

\begin{enumerate}
    \item \textbf{Языков высокого уровня.}

    Они позволяют выражать решения ближе к предметной области.

    \item \textbf{Структурного программирования.}

    Использование последовательности, ветвления и цикла вместо произвольных переходов облегчает понимание.

    \item \textbf{Объектно-ориентированного программирования.}

    Инкапсуляция, классы и интерфейсы помогают моделировать предметную область.

    \item \textbf{Функционального программирования.}

    Чистые функции и неизменяемые данные часто упрощают рассуждения о поведении программы.

    \item \textbf{Модульных систем.}

    Модули позволяют скрывать детали реализации.

    \item \textbf{Средств автодокументирования.}

    Документация может генерироваться из кода и комментариев.

    \item \textbf{Интеллектуальных сред разработки.}

    IDE помогают переходить к определениям, искать использования, строить диаграммы зависимостей.

    \item \textbf{Линтеров и форматтеров.}

    Они поддерживают единый стиль кода.

    \item \textbf{Средств рефакторинга.}

    Они позволяют улучшать структуру кода без изменения внешнего поведения.
\end{enumerate}

\subsection{Влияние требований познаваемости на технологию разработки}

Для обеспечения познаваемости в процессе разработки применяют:
\begin{itemize}
    \item архитектурное проектирование;
    \item код-ревью;
    \item соглашения о стиле;
    \item ограничение сложности функций;
    \item ведение технической документации;
    \item описание публичных интерфейсов;
    \item создание обучающих примеров;
    \item ведение журнала архитектурных решений;
    \item регулярный рефакторинг;
    \item парное программирование;
    \item коллективное владение кодом.
\end{itemize}

\section{Рациональная ресурсоемкость программного продукта}

\subsection{Определение ресурсоемкости}

\begin{definition}
\textbf{Ресурсоемкость программного продукта} --- это объем вычислительных, временных, сетевых, энергетических и иных ресурсов, необходимых программному продукту для выполнения своих функций.
\end{definition}

\begin{definition}
\textbf{Рациональная ресурсоемкость} --- это такое потребление ресурсов, при котором программный продукт удовлетворяет требованиям по времени отклика, пропускной способности, объему памяти, стоимости эксплуатации и масштабируемости без неоправданного перерасхода ресурсов.
\end{definition}

Рациональность не означает минимальное потребление ресурсов любой ценой. Иногда допустимо использовать больше памяти ради большей скорости или больше вычислений ради большей надежности.

\subsection{Виды ресурсов}

Программный продукт может потреблять:
\begin{itemize}
    \item процессорное время;
    \item оперативную память;
    \item дисковое пространство;
    \item сетевую пропускную способность;
    \item операции ввода-вывода;
    \item энергию;
    \item лицензии и платные облачные ресурсы;
    \item время пользователя;
    \item время администрирования.
\end{itemize}

\subsection{Основные показатели ресурсоемкости}

К показателям ресурсоемкости относятся:

\begin{enumerate}
    \item \textbf{Время отклика} --- время между запросом пользователя и ответом системы.
    \item \textbf{Пропускная способность} --- число операций, выполняемых за единицу времени.
    \item \textbf{Загрузка процессора}.
    \item \textbf{Потребление оперативной памяти}.
    \item \textbf{Количество операций ввода-вывода}.
    \item \textbf{Объем сетевого трафика}.
    \item \textbf{Время запуска}.
    \item \textbf{Время установки}.
    \item \textbf{Стоимость эксплуатации}.
    \item \textbf{Энергопотребление}.
\end{enumerate}

\subsection{Алгоритмическая эффективность}

Одним из главных факторов ресурсоемкости является выбор алгоритма.

\begin{definition}
\textbf{Временная сложность алгоритма} --- это зависимость количества элементарных операций от размера входных данных.
\end{definition}

\begin{definition}
\textbf{Пространственная сложность алгоритма} --- это зависимость объема дополнительной памяти от размера входных данных.
\end{definition}

Для оценки сложности используют асимптотические обозначения.

\begin{definition}
Функция $f(n)$ принадлежит классу $O(g(n))$, если существуют положительные константы $C$ и $n_0$ такие, что для всех $n \geq n_0$ выполняется:
\[
0 \leq f(n) \leq Cg(n).
\]
\end{definition}

Примеры:
\[
3n^2 + 10n + 5 = O(n^2),
\]
\[
7n + 100 = O(n),
\]
\[
\log n = O(n).
\]

\subsection{Пример влияния алгоритма на ресурсоемкость}

Задача: найти, есть ли в массиве два одинаковых элемента.

\subsubsection{Наивный алгоритм}

\textbf{Шаги алгоритма:}
\begin{enumerate}
    \item Взять первый элемент.
    \item Сравнить его со всеми следующими.
    \item Затем взять второй элемент.
    \item Сравнить его со всеми следующими.
    \item Продолжать до конца массива.
    \item Если найдена равная пара, вернуть \texttt{true}.
    \item Если пары нет, вернуть \texttt{false}.
\end{enumerate}

Псевдокод:
\begin{verbatim}
function has_duplicates(a):
    for i = 0 to n - 1:
        for j = i + 1 to n - 1:
            if a[i] == a[j]:
                return true
    return false
\end{verbatim}

Сложность:
\[
T(n) = O(n^2).
\]

Минимальный пример:
\[
a = [4, 7, 2, 7].
\]

Сравнения:
\[
4 \neq 7,\quad 4 \neq 2,\quad 4 \neq 7,\quad 7 \neq 2,\quad 7 = 7.
\]

Результат: найден дубликат.

\subsubsection{Алгоритм с множеством}

\textbf{Шаги алгоритма:}
\begin{enumerate}
    \item Создать пустое множество $S$.
    \item Последовательно просматривать элементы массива.
    \item Для очередного элемента $x$ проверить, есть ли он в $S$.
    \item Если $x \in S$, вернуть \texttt{true}.
    \item Иначе добавить $x$ в $S$.
    \item Если просмотр завершен, вернуть \texttt{false}.
\end{enumerate}

Псевдокод:
\begin{verbatim}
function has_duplicates(a):
    S = empty_set()
    for x in a:
        if x in S:
            return true
        add x to S
    return false
\end{verbatim}

Для хеш-таблицы средняя сложность:
\[
T(n) = O(n).
\]

Дополнительная память:
\[
M(n) = O(n).
\]

Иллюстрация для массива:
\[
a = [4, 7, 2, 7].
\]

\[
S = \varnothing.
\]

После обработки $4$:
\[
S = \{4\}.
\]

После обработки $7$:
\[
S = \{4,7\}.
\]

После обработки $2$:
\[
S = \{4,7,2\}.
\]

Следующий элемент $7$ уже принадлежит $S$, значит найден дубликат.

\subsubsection{Вывод}

Первый алгоритм требует меньше дополнительной памяти, но больше времени. Второй алгоритм быстрее в среднем случае, но требует дополнительной памяти. Это демонстрирует типичный компромисс:
\[
\text{время} \leftrightarrow \text{память}.
\]

\subsection{Методы снижения ресурсоемкости}

Основные методы:
\begin{enumerate}
    \item выбор более эффективных алгоритмов;
    \item выбор подходящих структур данных;
    \item кэширование;
    \item ленивые вычисления;
    \item пакетная обработка;
    \item асинхронное выполнение;
    \item параллелизм;
    \item уменьшение числа сетевых запросов;
    \item сжатие данных;
    \item оптимизация запросов к базе данных;
    \item профилирование;
    \item устранение утечек памяти;
    \item настройка параметров среды выполнения;
    \item горизонтальное и вертикальное масштабирование.
\end{enumerate}

\subsection{Кэширование как способ снижения ресурсоемкости}

Кэширование позволяет сохранять результат дорогой операции и повторно использовать его.

\textbf{Шаги алгоритма кэширования:}
\begin{enumerate}
    \item Получить запрос с ключом $k$.
    \item Проверить, есть ли результат в кэше.
    \item Если результат есть, вернуть его.
    \item Если результата нет, вычислить его.
    \item Сохранить результат в кэше.
    \item Вернуть результат.
\end{enumerate}

Псевдокод:
\begin{verbatim}
function get_data(k):
    if cache.contains(k):
        return cache[k]

    result = expensive_operation(k)
    cache[k] = result
    return result
\end{verbatim}

Минимальный пример:
\begin{itemize}
    \item запрос: курс валюты на дату;
    \item операция получения курса из внешнего сервиса занимает 500 мс;
    \item повторный ответ из кэша занимает 2 мс.
\end{itemize}

Кэширование снижает время отклика, но увеличивает потребление памяти и может приводить к проблеме устаревших данных.

\subsection{Влияние требований рациональной ресурсоемкости на системы программирования}

Эти требования влияют на развитие:

\begin{enumerate}
    \item \textbf{Оптимизирующих компиляторов.}

    Компиляторы выполняют оптимизации: удаление мертвого кода, разворачивание циклов, инлайнинг функций, оптимизацию регистров.

    \item \textbf{Профилировщиков.}

    Они показывают, где программа тратит время и память.

    \item \textbf{Сборщиков мусора.}

    Они управляют памятью, но должны делать это с приемлемыми накладными расходами.

    \item \textbf{Средств параллельного программирования.}

    Используются потоки, процессы, корутины, акторы, GPU-вычисления.

    \item \textbf{Средств асинхронного ввода-вывода.}

    Позволяют обслуживать много операций без блокировки потоков.

    \item \textbf{Специализированных структур данных.}

    Например, хеш-таблицы, B-деревья, очереди с приоритетом, фильтры Блума.

    \item \textbf{Средств мониторинга производительности.}

    Они фиксируют использование ресурсов в реальной эксплуатации.
\end{enumerate}

\subsection{Влияние требований рациональной ресурсоемкости на технологию разработки}

В процесс разработки включаются:
\begin{itemize}
    \item задание количественных требований к производительности;
    \item нагрузочное тестирование;
    \item стресс-тестирование;
    \item профилирование;
    \item анализ сложности алгоритмов;
    \item контроль потребления памяти;
    \item мониторинг после внедрения;
    \item оптимизация только после измерений;
    \item архитектурное проектирование с учетом масштабируемости.
\end{itemize}

\begin{principle}
\textbf{Принцип измеряемой оптимизации.}
Оптимизировать следует прежде всего те части программы, которые по результатам измерений действительно являются узкими местами.
\end{principle}

\section{Связь требований качества между собой}

Требования к качеству часто конфликтуют друг с другом. Поэтому задача проектирования состоит не в максимизации каждого свойства по отдельности, а в поиске разумного компромисса.

\begin{longtable}{p{0.25\textwidth} p{0.3\textwidth} p{0.35\textwidth}}
\toprule
\textbf{Требование} & \textbf{Может улучшать} & \textbf{Может ухудшать} \\
\midrule
Надежность & доверие к системе, устойчивость, безопасность данных & скорость разработки, простоту архитектуры, производительность из-за проверок и резервирования \\
\midrule
Переносимость & срок жизни продукта, независимость от платформы & производительность, использование специфичных возможностей платформы \\
\midrule
Познаваемость & сопровождаемость, качество разработки, скорость обучения & иногда краткость кода, иногда предельную производительность \\
\midrule
Рациональная ресурсоемкость & производительность, стоимость эксплуатации, масштабируемость & простоту реализации, переносимость, иногда надежность \\
\bottomrule
\end{longtable}

Пример компромисса:
\begin{itemize}
    \item использование низкоуровневого языка может повысить производительность;
    \item но может снизить безопасность памяти и повысить сложность сопровождения;
    \item использование высокоуровневого языка может повысить познаваемость и переносимость;
    \item но иногда увеличивает потребление памяти.
\end{itemize}

\section{Влияние требований к программному продукту на системы программирования}

\subsection{Понятие системы программирования}

\begin{definition}
\textbf{Система программирования} --- это совокупность языков, трансляторов, библиотек, сред разработки, отладчиков, средств сборки, тестирования и сопровождения, предназначенных для создания программ.
\end{definition}

В состав системы программирования могут входить:
\begin{itemize}
    \item язык программирования;
    \item компилятор или интерпретатор;
    \item стандартная библиотека;
    \item редактор или интегрированная среда разработки;
    \item отладчик;
    \item профилировщик;
    \item система сборки;
    \item пакетный менеджер;
    \item средства тестирования;
    \item статический анализатор;
    \item генераторы документации.
\end{itemize}

\subsection{Общее влияние требований качества}

Требования к программному продукту определяют, какие свойства должны иметь системы программирования.

\begin{longtable}{p{0.24\textwidth} p{0.66\textwidth}}
\toprule
\textbf{Требование} & \textbf{Требуемые свойства системы программирования} \\
\midrule
Надежность & строгая типизация, проверка границ, обработка исключений, анализаторы, тестовые фреймворки, отладчики, средства формальной проверки \\
\midrule
Переносимость & стандартизованный язык, кроссплатформенные библиотеки, виртуальная машина, независимая система сборки, пакетный менеджер \\
\midrule
Познаваемость & выразительный синтаксис, модульность, средства навигации по коду, автодокументация, форматтеры, линтеры, рефакторинг \\
\midrule
Рациональная ресурсоемкость & оптимизирующий компилятор, профилировщик, контроль памяти, средства параллелизма, эффективные структуры данных \\
\bottomrule
\end{longtable}

\section{Влияние требований на технологии разработки программных систем}

\subsection{Понятие технологии разработки}

\begin{definition}
\textbf{Технология разработки программных систем} --- это совокупность процессов, методов, инструментов, стандартов и организационных практик, применяемых для создания, проверки, внедрения и сопровождения программных систем.
\end{definition}

Технология разработки включает:
\begin{itemize}
    \item сбор и анализ требований;
    \item проектирование архитектуры;
    \item программирование;
    \item тестирование;
    \item сборку;
    \item развертывание;
    \item эксплуатацию;
    \item сопровождение;
    \item управление изменениями;
    \item управление качеством.
\end{itemize}

\subsection{Требования качества и этапы жизненного цикла}

\begin{longtable}{p{0.22\textwidth} p{0.68\textwidth}}
\toprule
\textbf{Этап} & \textbf{Влияние требований качества} \\
\midrule
Анализ требований & Требования должны быть измеримыми, проверяемыми и согласованными. Например, не ``система должна быть быстрой'', а ``95\% запросов должны выполняться не более чем за 200 мс''. \\
\midrule
Проектирование & Выбираются архитектурные стили, обеспечивающие надежность, переносимость, познаваемость и рациональное использование ресурсов. \\
\midrule
Реализация & Применяются стандарты кодирования, типизация, обработка ошибок, абстракции от платформы, эффективные алгоритмы. \\
\midrule
Тестирование & Проверяются не только функции, но и надежность, производительность, переносимость, удобство освоения. \\
\midrule
Развертывание & Используются автоматизированные сценарии установки, контейнеры, миграции, мониторинг. \\
\midrule
Эксплуатация & Собираются метрики отказов, времени отклика, потребления ресурсов, пользовательских ошибок. \\
\midrule
Сопровождение & Выполняются исправления, рефакторинг, перенос на новые платформы, оптимизация и обновление документации. \\
\bottomrule
\end{longtable}

\subsection{Архитектурные решения под требования качества}

\begin{longtable}{p{0.25\textwidth} p{0.65\textwidth}}
\toprule
\textbf{Требование} & \textbf{Типовые архитектурные решения} \\
\midrule
Надежность & резервирование, транзакции, очереди сообщений, контроль целостности, повторные попытки, автоматическое восстановление \\
\midrule
Переносимость & слоистая архитектура, изоляция платформо-зависимого кода, контейнеризация, стандартизованные интерфейсы \\
\midrule
Познаваемость & модульная архитектура, ясные границы компонентов, предметно-ориентированная модель, документация архитектурных решений \\
\midrule
Рациональная ресурсоемкость & кэширование, балансировка нагрузки, асинхронная обработка, масштабирование, оптимизация хранилищ данных \\
\bottomrule
\end{longtable}

\section{Формулирование качественных требований}

Качественные требования должны быть:
\begin{enumerate}
    \item \textbf{однозначными};
    \item \textbf{измеримыми};
    \item \textbf{проверяемыми};
    \item \textbf{реалистичными};
    \item \textbf{согласованными};
    \item \textbf{значимыми для заинтересованных сторон}.
\end{enumerate}

Плохие формулировки:
\begin{itemize}
    \item ``система должна быть надежной'';
    \item ``программа должна работать быстро'';
    \item ``интерфейс должен быть понятным'';
    \item ``продукт должен быть переносимым''.
\end{itemize}

Хорошие формулировки:
\begin{itemize}
    \item ``доступность сервиса должна быть не ниже $99.9\%$ в месяц'';
    \item ``95\% запросов поиска должны выполняться не более чем за 300 мс при нагрузке 100 запросов в секунду'';
    \item ``новый пользователь должен выполнить сценарий регистрации без обращения к справке не более чем за 3 минуты'';
    \item ``серверная часть должна запускаться в Linux и Windows без изменения исходного кода''.
\end{itemize}

\section{Итоговая схема влияния требований}

Требования к программному продукту оказывают влияние на все уровни создания программной системы:

\[
\text{Требования качества}
\rightarrow
\text{Архитектура}
\rightarrow
\text{Языки и инструменты}
\rightarrow
\text{Процесс разработки}
\rightarrow
\text{Тестирование}
\rightarrow
\text{Эксплуатация}.
\]

\begin{itemize}
    \item Надежность требует контроля ошибок, тестирования, резервирования и восстановления.
    \item Переносимость требует стандартизации, абстрагирования от платформы и проверки на разных средах.
    \item Познаваемость требует ясной структуры, документации, выразительных языков и дисциплины кодирования.
    \item Рациональная ресурсоемкость требует анализа алгоритмов, измерений, профилирования и оптимизации.
\end{itemize}

Главный вывод состоит в том, что требования к программному продукту нельзя рассматривать как второстепенные пожелания. Они определяют архитектуру, выбор языка программирования, инструменты, методы тестирования, организацию команды и стоимость жизненного цикла программной системы.

\end{document}
